
\chapter{PCI DSS}\label{ch:pci-dss}
\highlight{Архитектурный смысл PCI DSS сводится к~одному вопросу: где именно
в~системе появляются данные карты и~сколько компонентов должны их видеть}.
Ответ определяет, останется ли объём комплаенса в~отдельном сегменте или
распространится на~всю инфраструктуру.

\section{Двенадцать требований}\label{sec:pci-logic}

PCI DSS (Payment Card Industry Data Security Standard) определяет требования
к~защите данных платёжных карт для любой организации, которая хранит,
обрабатывает или передаёт данные держателя карты (cardholder data, CHD).
Стандарт разрабатывает и~поддерживает PCI SSC (Payment Card Industry Security
Standards Council), созданный глобальными карточными сетями~\cite{pci-dss-4}.

До появления PCI DSS у~карточных сетей действовали собственные программы защиты
данных. FAQ PCI SSC указывает, что общий стандарт возник из~выравнивания
программ Visa AIS/CISP (Account Information Security / Cardholder Information
Security Program) и~Mastercard SDP (Site Data
Protection)~\cite{pci-faq-1199}. Для ТСП, принимающего карты нескольких
брендов, это заменило набор частично несовпадающих программ общей базовой
планкой безопасности.

PCI SSC был создан в~2006 году American Express, Discover, JCB, Mastercard
и~Visa; разработанная учредителями версия 1.1 вступила в~силу вместе с~запуском
совета~\cite{pci-faq-1199,pci-faq-1227}. Утечки TJX и~Heartland показали масштаб
промышленной угрозы: компрометация затрагивает и~отдельные магазины, и~крупные
процессинговые компании~\cite{doj-tjx-2009,doj-heartland-2009}.

Карточные сети передавали требования Visa CISP и~Mastercard SDP через
эквайеров, а~не напрямую ТСП. Требования разных сетей расходились, поэтому
у~одного ТСП возникали несогласованные аудиты от~нескольких эквайеров. Прямого
финансового стимула инвестировать в~безопасность у~ТСП также не~было: убытки
от~компрометации распределялись через chargeback между эквайером и~эмитентом,
а~штрафы сетей выставлялись эквайеру. Единый стандарт изменил структуру
стимулов: несоблюдение PCI~DSS грозит ТСП штрафами и~повышенными ставками,
которые сеть взыскивает с~эквайера, а~эквайер взыскивает с~ТСП по~контракту.

Двенадцать требований PCI DSS группируются по~областям защиты. Требования~1--2
задают сетевую защиту: межсетевые экраны, сегментацию и~безопасные конфигурации
без заводских паролей и~лишних сервисов. Требования~3--4 относятся к~карточным
данным: требование~3 описывает защиту хранимого PAN и~запрет хранения Sensitive
Authentication Data (SAD), а~требование~4 предписывает шифровать передачу CHD
по~открытым сетям.

Требования~5--6 охватывают защиту систем от~вредоносного ПО, контроль
уязвимостей, разработку, патчи и~веб-приложения. Требования~7--9 сужают круг
людей и~систем, которые могут получить доступ к~карточным данным, и~требуют
физической и~логической дисциплины вокруг CDE (Cardholder Data Environment,
среда обработки карточных данных~--- см.~раздел~\ref{sec:pci-cde}).

Требования~10--11 требуют журналирования, мониторинга и~регулярного
тестирования безопасности среды, а~требование~12 закрепляет эти меры
в~политике информационной безопасности.

Если токенизация, P2PE или полностью вынесенная платёжная страница не допускают
данные карты в~среду ТСП, большая часть требований к~этой среде не применяется.

\section{CDE и~периметр}\label{sec:pci-cde}

\highlight{Определение периметра CDE~--- центральная задача в~работе по~PCI
DSS}. Ошибка в~определении периметра оценки означает либо недозащищённость,
когда данные карты оказываются вне CDE без должной защиты, либо избыточность,
когда организация тратит ресурсы на~системы, которые карточных данных никогда не
видят.

Scoping Information Supplement выделяет три типа систем в~периметре
оценки~\cite{pci-scoping-guidance}. Во-первых, это системы CDE, которые хранят,
обрабатывают или передают CHD/SAD: серверы авторизации, базы данных с~PAN,
терминалы. PCI DSS прямо запрещает хранить SAD (CVV2, full track, PIN/PIN block) после
авторизации даже в~зашифрованном виде; scoping не меняет это правило. Во-вторых, это непосредственно связанные системы
(connected-to systems): они могут не видеть карточные данные, но соединены
с~CDE и~способны повлиять на~него; типичный пример~--- DNS-сервер или иной
общий инфраструктурный узел в~том же сегменте. В-третьих, в~периметр попадают
security-impacting systems: серверы аутентификации, системы мониторинга,
антивирусные средства, серверы NTP и~прочие компоненты, от~которых зависит
безопасность самого CDE. Категории connected-to и~security-impacting относятся
к~области оценки, но не составляют сам CDE; Scoping Information Supplement
отдельно оговаривает, что строгого разделения между ними нет: одна и~та же
система может относиться к~обеим категориям.

\highlight{Сетевая сегментация~--- главный способ удержать CDE компактным}.
Если карточные данные находятся в~изолированном сегменте, отделённом
от~остальной инфраструктуры межсетевым экраном с~политикой запрета
по~умолчанию, в~периметр попадают этот сегмент и~его ближайшие зависимости.
Без такой изоляции CDE распространяется почти на~всю сеть организации,
и~\enquote{точечный комплаенс} теряет практический смысл.

Критерий сегментации задают правила межсетевого экрана: система, которая
не~может инициировать соединение с~CDE и~в~которую CDE не~может инициировать
соединение, остаётся вне периметра, даже если физически размещена в~той же
стойке. У~карточных данных нет сетевого пути в~неё даже при её полной
компрометации.

\practicenote{%
  ТСП выбрало SAQ~A (Self-Assessment Questionnaire~A,
  см.~раздел~\ref{sec:pci-saq}), но разработчик добавил собственный скрипт
  на~платёжную страницу. Сценарий перестаёт соответствовать SAQ~A и~переходит
  в~SAQ~A-EP
  (Eligible Merchants with Partially Outsourced E-commerce Payment Channels)
  или~SAQ~D, потому что скрипт ТСП участвует в~обработке данных карты.
  QSA (Qualified Security Assessor, аккредитованный аудитор PCI)
  фиксирует это на~аудите; после этого следуют штрафы и~срочная переработка
  интеграции. Изменение кода платёжной страницы меняет периметр PCI.

}

\begin{figure}[tbp]
  \centering
  \includefiguresvg[width=\linewidth]{assets/figures/ch12-pci-dss/pci-cde-boundary}
  \caption{Периметр CDE и~смежные системы.}\label{fig:pci-cde-boundary}
\end{figure}

\section{Сужение периметра: токенизация, P2PE и~вынесенные формы}\label{sec:pci-scope-reduction}

Периметр оценки сужается только там, где архитектура сокращает число систем,
способных увидеть PAN или повлиять на~среды, в~которых PAN существует.

\subsection*{Токенизация}

В~контексте PCI токенизация имеет три разных варианта, и~их нельзя смешивать;
терминологическая матрица \enquote{термин~$\times$~источник} и~семейство
токенов по~области действия разобраны
в~разделах~\ref{sec:token-terminology} и~\ref{sec:token-scope}
главы~\ref{ch:tokenization}. Первый вариант~--- сетевые токены, описанные
в~разделе~\ref{sec:token-anatomy}; они включают DPAN и~MPAN, которые выпускает
карточная сеть. Сетевые токены снижают ценность утечки и~позволяют перестроить
архитектуру так, чтобы настоящий PAN не проходил через системы ТСП. Подмена PAN
токеном не выводит систему из~периметра PCI сама по~себе. Решающие вопросы:
может ли среда увидеть PAN, использовать токен как полноценный платёжный
идентификатор или повлиять на~безопасность систем, где PAN
существует~\cite{pci-dss-4}.

Второй вариант~--- токенизация на~стороне PSP (Payment Service Provider,
платёжный провайдер). ТСП один раз передаёт PAN в~платёжный шлюз, шлюз хранит
его в~собственном хранилище токенов (vault) и~возвращает суррогатное значение
(payment token). После этого ТСП работает только с~этим суррогатом. Vault переносит основную
нагрузку на~PSP, но у~ТСП остаётся периметр вокруг точки первичного ввода PAN
и~интеграции со~шлюзом. Шлюз, владеющий vault, остаётся полноценным
субъектом PCI DSS, а~точка первичного ввода PAN требует отдельной
защиты~\cite{pci-dss-4}.

Третий вариант~--- vault или FPE на~стороне самой организации, без участия PSP
или карточной сети. Vault при первом приёме PAN генерирует случайный суррогат
(в~терминах PCI~SSC~--- \enquote{surrogate value}, в~отраслевой терминологии~---
\enquote{pseudo PAN}). Синтаксически он валиден как PAN, но
не~маршрутизируется в~сети и~к~авторизации
непригоден~\cite{pci-tokenization-guidelines}. Соответствие хранится
в~защищённом хранилище. Нижестоящие (\textit{downstream}) системы~--- аналитика,
биллинг, CRM~--- работают только с~суррогатом и~не видят PAN. В~FPE
(Format-Preserving Encryption, алгоритм FF1 по~NIST SP~800-38G)\footnote{Для
  FF1 NIST рекомендует ограничивать применение доменами, в~которых число
  возможных входов ($\mathrm{radix}^{\mathrm{minlen}}$) не меньше $10^6$, поэтому
  на~одиночные короткие поля он не применяется. Соседний алгоритм FF3 после атаки
  Durak--Vaudenay 2017~года был заменён на~FF3-1 в~черновике SP~800-38G
  Revision~1, а~во второй редакции черновика (2025) предложен к~полному
исключению.} PAN шифруется в~значение той же длины и~символьного состава без
отдельного хранилища соответствий; детокенизация требует ключа. В~обоих случаях
CDE сужается до~vault или инфраструктуры ключей и~точки первичного попадания
PAN в~систему. Прочие системы выходят из~периметра. Ни карточная сеть, ни НСПК,
ни внешний API здесь не
задействованы~\cite{pci-dss-4,pci-tokenization-guidelines,nist-sp-800-38g}.

Для российских кредитных организаций, работающих по~Положению Банка России
от~30.01.2025~№~851-П и~ГОСТ~Р~57580.1-2017 (национальному стандарту по~защите
информации финансовых организаций)~\cite{cbr-851p,gost-57580-1}, выбор между
vault и~FPE неравнозначен. Аналога FPE по~ГОСТ не существует~--- ТК 26
(Технический комитет
  по~стандартизации \enquote{Криптографическая защита информации} при
Росстандарте) не стандартизировал ни один режим
FPE~\cite{gost-34-13,tc26-recommendations}, ни один российский производитель не
имеет сертификата ФСБ (Федеральной службы безопасности) на~конструкцию FPE~---
ни на~базе Кузнечика, ни на~базе AES. В~vault сам токен не несёт
криптографического содержания и~не подпадает под требования к~СКЗИ (средствам
криптографической защиты информации в~терминологии ФСБ). Требования
к~сертифицированной криптографии применяются к~хранилищу соответствий
(шифрование хранимых PAN~--- Кузнечик через сертифицированный СКЗИ)
и~к~транспорту (ГОСТ~TLS при обращении к~vault по~сети)~\cite{cbr-851p}.

\subsection*{P2PE (Point-to-Point Encryption)}

P2PE~--- стандарт PCI, при котором данные карты шифруются непосредственно
в~терминале в~момент считывания и~расшифровываются только в~HSM на~стороне
эквайера или процессора. Центральный элемент решения~--- модуль SRED (Secure
Reading and Exchange of Data) внутри терминала. Это аппаратно защищённый
компонент, который перехватывает считанные данные до их попадания в~общую
память терминала
и~шифрует их на~ключе, введённом во время процедуры key injection.
В~результате через инфраструктуру ТСП проходит только зашифрованный блок данных,
а~не открытый PAN~\cite{pci-p2pe}.

Модель P2PE разделяет ответственность между двумя ролями: P2PE Component
Provider, то есть поставщиком отдельных компонентов (терминала, HSM, системы
управления ключами), и~P2PE Solution Provider, то есть поставщиком
сертифицированного конечного решения. ТСП с~готовым решением P2PE получает
меньше применимых требований и~может валидироваться по~более узкому сценарию,
если это допускают условия его платёжной программы и~требования бренда карты.
Точкой расшифровки остаётся HSM процессора, и~эта среда сохраняет полный
периметр PCI DSS и~несёт ответственность за~защиту
ключей~\cite{pci-p2pe}.

\subsection*{Вынесенные поля оплаты и~iFrame}

В~электронной коммерции ТСП может не собирать данные карты на~собственной
странице. Вместо этого он встраивает iFrame или вынесенные поля оплаты (hosted
fields), которые предоставляет PSP. Ответственность здесь распределяется
по~принадлежности элемента DOM. Поле ввода размещено в~домене PSP, нажатия
клавиш отправляются прямо туда, и~ни PAN, ни CVV не касаются ни кода JavaScript
ТСП, ни его серверов. От~PSP ТСП получает платёжный токен и~работает только
с~ним~\cite{pci-dss-4}.

Если \emph{все} элементы платёжной страницы поставляются только провайдером,
совместимым с~PCI DSS, ТСП может документировать соответствие через SAQ~A. Если
хотя бы часть элементов формируется на~стороне ТСП, применяется SAQ~A-EP.
SAQ~A v4.0.1~(r1), действующий с~января 2025~года, требует для сценария
со~встроенной платёжной формой (iframe или скрипт PSP) подтвердить, что весь
сайт ТСП защищён от~атак через скрипты, способных скомпрометировать ввод данных
карты. Для чистого redirect и~полностью вынесенной оплаты PCI SSC в~FAQ~1588
это условие не
распространяет~\cite{pci-saq-a-r1-2025,pci-faq-1588,pci-faq-1291,pci-faq-1293}.

\subsection*{Полное перенаправление на~страницу PSP}

Полное перенаправление (redirect) остаётся наиболее жёстким архитектурным
способом вывести среду ТСП из~CDE. Во время оплаты держатель карты целиком покидает
сайт ТСП. Браузер перенаправляется на~страницу платёжного шлюза, карточные
данные вводятся и~обрабатываются \emph{целиком на~стороне PSP}. После
завершения держатель возвращается с~уведомлением об~успехе или отказе. PAN, CVV
и~срок действия карты в~среде ТСП не появляются, и~при выполнении условий
вопросника этот путь ведёт к~SAQ~A~\cite{pci-saq}.

Покупатель при этом замечает смену домена в~адресной строке, и~часть таких
сессий заканчивается брошенной корзиной. PSP смягчают эффект стилизованными
страницами оплаты, визуально близкими к~сайту ТСП. Для малого бизнеса и~для
случаев, где упрощение комплаенса важнее непрерывности интерфейса, redirect
остаётся предпочтительным выбором~\cite{pci-saq}.

Токенизация (как сетевая, так и~vault/FPE на~стороне ТСП или PSP) остаётся
самостоятельным путём сужения периметра и~не сводится к~выбору канала ввода:
даже при сохранённой собственной форме оплаты вынос PAN из~нижестоящих систем
в~vault или замена PAN на~DPAN сокращает число систем, остающихся в~CDE.
При сужении периметра ответственность за~CHD переносится на~PSP или
валидированное P2PE-решение; процесс всё равно остаётся подконтрольным PCI DSS.

\section{SAQ~A, SAQ~A-EP и~SAQ~D}\label{sec:pci-saq}

SAQ (Self-Assessment Questionnaire)~--- шаблон самооценки, который PCI~SSC
публикует; ТСП заполняет его и~передаёт эквайеру по~программе соответствия
карточной сети без полного аудиторского цикла с~участием QSA. Тип SAQ
определяется конкретным способом взаимодействия с~карточными данными~--- тем,
как карта попадает в~платёжный процесс и~соприкасается ли ТСП
с~CHD~\cite{pci-saq}.

Интернет-магазин, который при~оплате полностью перенаправляет покупателя
на~страницу PSP по~сценарию redirect, использует \textbf{SAQ~A}~--- наиболее
узкий вопросник для типичного ТСП в~электронной коммерции с~вынесенной
страницей оплаты. Но SAQ~A не означает нулевой периметр. PCI SSC отдельно
уточняет, что веб-сервер ТСП, на~котором расположен механизм redirect, остаётся
в~периметре, чтобы аудитор мог проверить конфигурацию и~сам механизм
перенаправления~\cite{pci-faq-1332}.

Со~встроенным iFrame критерий строже: если все элементы платёжной страницы,
которые получает браузер, поступают только напрямую от~поставщика услуг,
совместимого с~PCI DSS, реализация может оставаться в~SAQ~A. Если же хотя бы
часть элементов страницы формируется сайтом ТСП, применяется \textbf{SAQ~A-EP}.
Так PCI SSC различает Direct Post и~iframe/URL redirect: в~Direct Post сайт ТСП
сам участвует в~приёме карточных данных, тогда как в~iframe и~redirect страницу
для ввода реквизитов формирует процессор~\cite{pci-faq-1291,pci-faq-1293}. Для
встроенной формы в~SAQ~A~r1 добавлено условие~--- ТСП должно подтвердить, что
его страница не подвержена атакам через скрипты; для чистого redirect
и~полностью вынесенной оплаты это условие не
применяется~\cite{pci-saq-a-r1-2025,pci-faq-1588}.

Для сценариев с~физическими терминалами действуют отдельные типы вопросников.
\textbf{SAQ~B-IP} применяется к~торговым точкам с~терминалами,
сертифицированными по~PTS (PIN Transaction Security~--- программа PCI~SSC
по~сертификации устройств приёма карт), подключёнными к~платёжному процессору
по~IP. Для PCI DSS~v4 PCI SSC отдельно уточнил, что этот тип рассчитан
на~автономные point-of-interaction (POI) устройства, одобренные PCI SSC и~не
смешанные с~другими типами устройств в~той же сетевой зоне~\cite{pci-faq-1443}.

\textbf{SAQ~C-VT} покрывает другой класс сценариев~--- колл-центр или бэк-офис,
где сотрудник вручную вводит карточные данные через виртуальный терминал PSP,
то есть через веб-интерфейс, который сам не хранит CHD. Официальное FAQ PCI SSC
дополнительно оговаривает: это изолированный компьютер с~браузером, через
который вводится по~одной транзакции за~раз~\cite{pci-faq-1229}.

Отдельные типы: \textbf{SAQ~B} для торговых точек с~импринтерами или
автономными терминалами с~коммутируемым (\textit{dial-out}) соединением, без
хранения CHD в~электронном виде; \textbf{SAQ~P2PE}~--- наиболее узкий вопросник
для ТСП, использующих валидированное PCI~P2PE-решение без доступа к~открытым
карточным данным.

Если архитектуру нельзя свести ни к~одному из~упрощённых сценариев, остаётся
\textbf{SAQ~D}. Он применяется ко~всем ТСП, которые сами хранят, обрабатывают
или передают CHD на~собственных серверах (не~только хранение, но и~любая
обработка или транзит) и~выходят за~пределы специальных типов, а~также ко~всем
поставщикам услуг без исключения. SAQ~D соответствует полному объёму требований
PCI DSS. Разница между SAQ~A и~SAQ~D~--- в~трудоёмкости: больше тестов, больше
документации, ежеквартальные сканы ASV (Approved Scanning Vendor,
авторизованный поставщик внешнего сканирования уязвимостей), ежегодный pentest
(penetration test, тест на~проникновение). Переход с~SAQ~D на~SAQ~A для
большинства ТСП окупается в~первый комплаенс-цикл~\cite{pci-saq}.

Для крупных ТСП (Level~1~--- уровень, который карточная сеть назначает
  мерчантам с~годовым объёмом свыше~6\,млн транзакций по~своему бренду или после
крупного инцидента) и~поставщиков услуг (service providers~--- организации,
обрабатывающие, хранящие или передающие CHD от~имени других) SAQ~D дополняется
полным аудитом QSA с~выпуском ROC (Report on Compliance, отчёт о~соответствии
PCI). Конкретная стоимость соответствия зависит от~размера CDE, сложности
инфраструктуры, региональных тарифов QSA и~состава внешних сервисов; публичных
усреднённых тарифов PCI SSC не публикует, поэтому числовые ориентиры в~книге не
приводятся~\cite{pci-dss-4}.

\begin{figure}[tbp]
  \centering
  \includefiguresvg[width=\linewidth]{assets/figures/ch12-pci-dss/pci-saq-decision}
  \caption{Восемь типов SAQ и~условия их применимости.}\label{fig:pci-saq-decision}
\end{figure}

\section{PCI DSS~4.0.1}\label{sec:pci-401}

Версия~4.0.1 ужесточила контроль платёжной страницы, инвентаря скриптов
и~доказуемости внутренних процедур. Ключевые требования направлены против атак
семейства Magecart, при которых злоумышленник внедряет вредоносный JavaScript
на~платёжную страницу для перехвата карточных данных~\cite{pci-dss-4}.

\textbf{Требование~6.4.3} обязывает инвентаризировать, обосновывать
и~авторизовывать каждый JavaScript, загружаемый на~странице ввода карточных
данных. Для интернет-торговли это увеличивает объём работ: типичная платёжная
страница подгружает десятки скриптов~--- аналитика, A/B-тесты, менеджеры тегов,
виджеты чата~--- и~каждый требует документирования и~контроля целостности
(CSP~--- Content Security Policy; Subresource Integrity, SRI~--- механизмы W3C
для контроля загружаемых браузером скриптов).

\textbf{Требование~11.6.1} обязывает обнаруживать несанкционированные изменения
HTTP-заголовков и~содержимого платёжной страницы: контроль целостности,
клиентский мониторинг и~процедуры, рассчитанные на~электронный скимминг.

Многофакторная аутентификация (MFA, Multi-Factor Authentication) стала
обязательной для всех обращений к~CDE, а~не только для удалённого доступа, как
в~3.2.1. Внутренние сканы уязвимостей должны выполняться
в~аутентифицированном режиме, чтобы выявлять конфигурационные и~локальные
проблемы, которые не видны снаружи. Для защиты PAN при передаче стандарт
строже фиксирует работу с~доверенными ключами и~сертификатами: их нужно
инвентаризировать, принимать только доверенные экземпляры и~контролировать
валидность. Targeted Risk Analysis требует документировать периодичность
контрольных мер, если стандарт допускает гибкость.

Большинство новых требований версии~4.0 стали обязательными после 31~марта
2025~года, когда завершился переходный период, в~течение которого их
допускалось трактовать как будущие обязательные практики~\cite{pci-dss-4}.

\section{Customized Approach}\label{sec:pci-customized}

PCI DSS~4.0 ввёл альтернативный путь выполнения требований~--- Customized
Approach (индивидуальный путь с~собственными контролями и~доказательством их
  эффективности; противопоставляется Defined Approach~--- стандартному выполнению
требований по~перечню обязательных контролей). Организация вправе достичь той
же цели безопасности альтернативным контролем при условии, что докажет
эквивалентность на~проверке~\cite{pci-dss-4}.

Для использования Customized Approach организации нужно сначала зафиксировать
цель безопасности, которую стандарт связывает с~конкретным требованием, затем
провести целевой анализ риска (Targeted Risk Analysis), обосновать, что
альтернативный контроль достигает той же цели, и~предоставить аудитору QSA
документацию и~доказательства, достаточные для валидации.

TRA по~Customized Approach начинается с~Customized Approach Objective~--- цели
безопасности для конкретного требования, как её определяет стандарт. Затем
описывается текущее окружение, угрозы и~уязвимости, относящиеся к~этой цели.
После этого документируется реализованный контроль: какие выявленные угрозы он
нейтрализует и~почему достигает заявленной цели. Остаточный риск фиксируется
в~документации, свидетельства подготавливаются для QSA. Документация должна
быть достаточной для независимой валидации; устная аргументация не заменяет
записи~\cite{pci-dss-4}.

Например, организация строит микросервисную архитектуру с~zero-trust
сегментацией (модель, при которой ни одна сеть или узел не считается доверенным
по~умолчанию, а~доступ выдаётся по~идентичности и~контексту) вместо
классических межсетевых экранов. Требование~1 (Install and Maintain Network
Security Controls) в~Defined Approach предполагает традиционные firewall rules.
В~Customized Approach организация обосновывает, что mTLS (mutual TLS~---
взаимная аутентификация сторон по~X.509-сертификатам) между сервисами,
identity-aware proxy и~непрерывный мониторинг сетевого трафика достигают той же
цели изоляции CDE, но иным способом. QSA должен независимо оценить
эффективность контроля.

TRA требует зрелой программы ИБ с~формализованным управлением рисками,
описанными контролями и~регулярным пересмотром. QSA тратит больше времени
на~валидацию нестандартных контролей, чем на~проверку по~списку Defined
Approach. Для организации это означает и~более высокие консалтинговые расходы,
и~более длительный аудит.

Для типовой архитектуры Defined Approach проще и~дешевле. Customized Approach
оправдан только там, где архитектура не укладывается в~стандартный шаблон, но
безопасность можно убедительно доказать.

\section{PCI DSS и~российские требования: 821-П, 851-П}\label{sec:pci-russian}

Для организации, работающей с~карточными платежами в~России, PCI DSS никогда не
действует в~одиночку. Поверх требований карточных сетей действуют Положение
Банка России от~17.08.2023~№~821-П \enquote{О~защите информации при~переводах
денежных средств} и~Положение Банка России от~30.01.2025~№~851-П
\enquote{О~противодействии переводам без согласия
клиента}~\cite{cbr-821p,cbr-851p}~(см.~раздел~\ref{sec:regulation-cbr-provisions});
конкретный состав мер они задают через нормативную отсылку
к~ГОСТ~Р~57580.1-2017, а~оценку соответствия~--- по~ГОСТ~Р~57580.2-2018%
\cite{gost-57580-1,gost-57580-2}~(разбор стандарта
в~разделе~\ref{sec:regulation-57580}).

И~PCI DSS, и~российские положения требуют криптографической защиты данных,
разграничения доступа, идентификации пользователей, MFA, журналирования,
мониторинга, реагирования на~инциденты и~управления уязвимостями. Разница
в~том, что российские требования дополнительно опираются на~сертифицированные
СКЗИ и~алгоритмы ГОСТ там, где это предписано регулятором
(раздел~\ref{sec:crypto-gost}), тогда как PCI DSS не навязывает конкретный
алгоритм.

Обязательность алгоритмов ГОСТ опирается на~сертификацию ФСБ~---
криптографический продукт должен пройти российскую процедуру оценки.
Иностранные реализации, включая AES-256 и~RSA, этой сертификации не имеют,
поэтому там, где регулятор требует ГОСТ, они применяться не могут.

PCI DSS допускает AES, RSA и~иную \enquote{strong cryptography} (в~терминологии
  PCI~--- алгоритмы и~длины ключей, удовлетворяющие действующим рекомендациям
NIST и~ISO), тогда как 821-П и~851-П для указанных регулятором сценариев
требуют российские криптографические механизмы и~сертифицированные средства
защиты~\cite{cbr-821p,cbr-851p}. PCI DSS предполагает ROC, AOC (Attestation of
Compliance, аттестат соответствия) или SAQ перед карточной сетью или эквайером,
а~российские требования добавляют собственную оценку соответствия, регуляторную
отчётность и~надзор Банка России. 851-П требует для~категорий прикладного ПО,
перечисленных в~положении, либо сертификации, либо оценки соответствия не ниже
ОУД4 (оценочный уровень доверия~4 по~ГОСТ~Р~ИСО/МЭК~15408); для части иного
прикладного ПО необходимость такой оценки банк определяет сам~\cite{cbr-851p}.

Прямого аналога этому в~PCI DSS нет. Требование~6 PCI DSS описывает внутренние
процессы безопасной разработки, управление уязвимостями и~жизненный цикл
изменений, но не вводит сопоставимой внешней процедуры сертификации или
формальной оценки соответствия для~прикладного банковского ПО. Для классов
банковских приложений, попадающих в~область применения 851-П, набор требований
шире PCI DSS~\cite{cbr-851p}.

Практичный подход к~параллельному комплаенсу~--- общая матрица контрольных мер,
в~которой каждому требованию PCI DSS сопоставлены пункты 821-П и~851-П. Матрица
делится на~три блока:

\begin{itemize}
  \item \textbf{пересечения}~--- меры, удовлетворяющие оба свода одновременно:
    контроль доступа, журналирование, сетевая сегментация, управление
    инцидентами; большая часть требований;
  \item \textbf{технические расхождения}~--- криптооперации внутри периметра
    реализуются на~алгоритмах ГОСТ с~сертифицированными СКЗИ, тогда как тот
    же трафик дополнительно защищается TLS для~требования~4;
  \item \textbf{непересечения}~--- сертификация или оценка соответствия
    прикладного ПО (только 851-П), регуляторная отчётность и~контроль Банка
    России (только 821-П/851-П), SAQ/ROC перед карточной сетью (только PCI
    DSS).
\end{itemize}

Объединённая программа с~общей доказательной базой и~QSA, знающим российское
регулирование, сокращает дублирующий сбор свидетельств по~двум сводам
и~уменьшает общую трудоёмкость аудита.

\highlight{PCI DSS и~57580 не~заменяют друг друга: для~российской финансовой
организации в~карточном периметре обязательны оба свода одновременно}.
Разделение проходит по~объекту защиты, а~не по~юрисдикции
аудитора. PCI DSS обращён к~данным платёжных карт (CHD/SAD) у~любого участника
карточной сети; 57580 обращён к~информации финансовой организации в~целом,
и~карточный периметр в~нём~--- лишь одна из~зон применения.

Там, где обе области пересекаются~--- эквайринг, эмиссия, процессинг внутри
лицензированного периметра~--- организация одновременно проходит оценку
соответствия по~57580.2 (через проверяющую организацию с~лицензией ФСТЭК
и~надзор Банка России) и~аттестацию по~PCI DSS (через QSA и~отчёт перед
карточной сетью).

Две оценки проходят разными циклами с~разной длительностью и~разной точкой
отсчёта. PCI DSS требует ежегодного подтверждения соответствия: SAQ вместе
с~AOC или ROC вместе с~AOC завершаются раз в~год, и~12-месячный цикл
отсчитывается от~даты завершения предыдущей оценки (Date of Report для~ROC,
Self-Assessment completion date для~SAQ). Эквайеру и~карточной сети передаётся
AOC; ROC остаётся у~QSA и~ассессируемой стороны
и~предоставляется только по~запросу~\cite{pci-saq}. Оценка соответствия
по~57580.2 проводится не~реже одного раза в~два года: периодичность задаёт
Положение Банка России от~30.01.2025~№~851-П \enquote{О~противодействии
переводам без согласия клиента} (п.~13.4)~\cite{cbr-851p}, а~саму методику~---
ГОСТ~Р~57580.2-2018, утверждённый приказом Росстандарта
от~28.03.2018~№~156-ст~\cite{gost-57580-2}. Период отсчитывается от~даты
завершения предыдущей внешней оценки, отчёт хранится не~менее пяти лет (851-П,
п.~13.2; см.~раздел~\ref{sec:regulation-cbr-provisions}). Внеплановые контрольные
мероприятия Банка России возможны вне этого расписания на~основании ст.~73
и~73.1 Федерального закона от~10.07.2002~№~86-ФЗ \enquote{О~Центральном банке
Российской Федерации}~\cite{fz-86}. QSA и~проверяющая организация ФСТЭК
работают независимо и~не~обмениваются свидетельствами; общая часть
доказательной базы~--- конфигурации, журналы, политики, результаты сканов~---
собирается организацией один раз, но интерпретируется под два свода отдельно.
Если контроль реализован альтернативным способом и~закрывает требование PCI DSS
через Customized Approach, российский регулятор всё равно требует обоснования
по~57580. ГОСТ~Р~57580.1-2017 предусматривает \enquote{иные (компенсирующие)
меры} (п.~6.4): при~невозможности технической реализации выбранной меры
финансовая организация разрабатывает компенсирующую меру, направленную
на~нейтрализацию тех же угроз из~модели угроз~\cite{gost-57580-1}.
Это~--- аналог Compensating Controls в~Defined Approach, а~не Customized
Approach: симметрии с~Targeted Risk Analysis и~отдельным Customized Approach
Objective в~57580 нет.
